\clearpage
\section{Anexo de reprodutibilidade e rastreabilidade}
\label{app:reproducibility}

O projeto LaTeX é autocontido: as entradas das figuras, tabelas,
listagens e bibliografia são locais. \path{dados/benchmark.csv}
e \path{dados/benchmark-runs.csv} são cópias literais dos CSVs
analisados; \path{dados/instancias/} contém os grafos e o
manifesto. Os seis PNGs em \path{figures/} são os originais
do benchmark. Os hashes completos e a auditoria constam de
\path{evidencias/}.

O relatório usa os arquivos \path{sbc-template.sty} e
\path{sbc.bst} sem alterações, obtidos do template SBC
fornecido pelo grupo e identificado na página oficial do
Overleaf.\footnote{\raggedright\url{https://www.overleaf.com/latex/templates/sbc-conferences-template/blbxwjwzdngr}}
O documento preserva papel A4, coluna única, margens,
fonte principal, cabeçalho de autores, legendas,
ausência de numeração de páginas e estilo bibliográfico
do modelo. O material inclui uma cópia do template de
origem para conferência.

\subsection{Reproduzir a aplicação e o protocolo}

Na raiz do projeto Java, os comandos a seguir compilam e
testam a aplicação. A geração usa as mesmas sementes;
uma nova medição pode produzir tempos, timeouts e
incumbentes diferentes, e não deve substituir
silenciosamente os dados auditados.

\begin{lstlisting}[language={},numbers=none,basicstyle=\ttfamily\footnotesize,caption={Compilação e testes usando somente o JDK 17.},label={lst:commands}]
javac -encoding UTF-8 --release 17 -Xlint:all -d out src/*.java tests/*.java
java "-Dfile.encoding=UTF-8" "-Djava.awt.headless=true" -cp out TestRunner
java -cp out Main generate --output data/instances/ --seed 20261115
\end{lstlisting}

O comando completo da medição original é mostrado na
Listagem~\ref{lst:benchmark-command}. Deve ser informado
em uma única linha; as quebras exibidas na listagem são
apenas visuais. Para guardar uma nova execução, escolha
outro caminho de saída.

\begin{lstlisting}[language={},numbers=none,basicstyle=\ttfamily\footnotesize,caption={Configuração do benchmark original, com três medições por algoritmo e instância.},label={lst:benchmark-command}]
java -Xms128m -Xmx512m "-Dfile.encoding=UTF-8" "-Djava.awt.headless=true" -cp out Main benchmark --timeout-ms 5000 --repetitions 3 --warmup 100 --output data/results/benchmark.csv
\end{lstlisting}

Para uso didático, o \path{executar.cmd} da aplicação
compila e abre um menu; seu experimento rápido de
oito grafos não é a amostra deste artigo.
A interface completa mantém os comandos individuais.

\subsection{Reproduzir o pós-processamento e o PDF}

Na raiz deste projeto LaTeX, o utilitário
\path{ferramentas/ReportStats.java} audita os dados e
regenera as tabelas, sem medir tempos novos:

\begin{lstlisting}[language={},numbers=none,basicstyle=\ttfamily\footnotesize,caption={Pós-processamento separado da aplicação experimental.},label={lst:postprocess}]
javac -encoding UTF-8 --release 17 -d build/java codigo/src/*.java codigo/tests/*.java ferramentas/ReportStats.java
java "-Dfile.encoding=UTF-8" "-Djava.awt.headless=true" -cp build/java ReportStats
\end{lstlisting}

O \texttt{main.tex} compila com pdfLaTeX e BibTeX.
No Overleaf, basta importar o ZIP fornecido e selecionar
\path{main.tex} como documento principal, usando
pdfLaTeX. O Java não é necessário para compilar o PDF,
pois as tabelas e figuras já estão incluídas. A
Listagem~\ref{lst:latex-build} mostra a sequência local.

\begin{lstlisting}[language={},numbers=none,basicstyle=\ttfamily\footnotesize,caption={Compilação local do artigo e resolução de referências cruzadas.},label={lst:latex-build}]
pdflatex -interaction=nonstopmode -halt-on-error main.tex
bibtex main
pdflatex -interaction=nonstopmode -halt-on-error main.tex
pdflatex -interaction=nonstopmode -halt-on-error main.tex
\end{lstlisting}

Os nomes completos dos autores constam dos metadados do
PDF e do README; a linha do título utiliza iniciais
intermediárias para manter os três autores em uma linha.
O pacote eletrônico permite conferir os dados e os
fontes independentemente da publicação posterior em um
repositório ou projeto compartilhado.
